\documentclass[book.tex]{subfiles}
\begin{document}

\chapter{Geometric Operators for Shape Analsis}
\label{chap:geometric_shape}
\noindent\rule{11cm}{0.4pt} \\
\textbf{Prerequisites:}  \autoref{chap:molecular_hamiltonian},   \\
\textbf{Difficulty Level:} ** \\
\noindent\rule{11cm}{0.4pt}
\vspace{5mm}

"""
Given many variants of geometry representations (meshes with different samplings, point clouds, triangle soups), the operator approach has been adopted to design algorithms that are invariant to the geometry representation. The operator approach considers the discretization of some continuous operator, whose behavior mimics that of the continuous counterpart, thus remaining stable and consistent across geometry representations. While the intrinsic Laplace–Beltrami operator has been a ubiquitous choice thanks to its often desirable properties such as invariance to isometries, it fails to capture the spatial embedding of a shape because it discards extrinsic information. To address these challenges, several alternative operators for shape analysis have been proposed in recent work, including the Dirichlet-to-Neumann operator and the Dirac operator, in addition to variants of modified Laplacian operators. I will give an extensive survey on these operators and explain how these operators can be utilized in the geometric learning setting. Finally, I will talk about a neural architecture to learn task-specific geometric operators, whose design is inspired by the success of the operator approach to shape analysis.
"""

References \cite{wang2018steklov}, \cite{wang2019intrinsic}, \cite{wang2019learning}.

\end{document}